\paragraph{温度参数延拓与 Bernoulli 子图矩：一般图配分函数、Hankel 正性与例外块端点刚性}
在软核复制链口径中，传递矩阵 \(T^{(q)}\) 对应于“重叠边”代价 \(1/2\) 的特例。
本段将该代价推广为温度门参数 \(p\in[0,1]\)，并把环幂迹与一般图上的配分函数严格识别为 Bernoulli 随机子图之独立集计数的 \(q\)-次矩。
该识别一方面在 \(q\)-轴上自动给出 Stieltjes 矩结构与 Hankel 正半定，另一方面在 \(m\)-轴上诱导一个带权的二字母词迹展开，并保持“固定谱簇幂和的单项塌缩”机制在全参数族内的稳定性。

\begin{definition}[温度门软核与 Bernoulli 子图矩]\label{def:pom-replica-softcore-temperature-kernel-bernoulli}
取整数 \(q\ge 1\)，令 \(U_q:=\{0,1\}^q\)，并以 \(u\wedge v\) 表示逐坐标与。
对 \(p\in[0,1]\) 定义温度门软核 \(T^{(q)}_p\in\RR^{U_q\times U_q}\)：
\begin{equation}\label{eq:pom-replica-softcore-temperature-kernel}
\boxed{\;
T^{(q)}_p(u,v)
:=
\begin{cases}
1,& u\wedge v=0,\\[2pt]
p,& u\wedge v\neq 0,
\end{cases}
\qquad(u,v\in U_q).\;}
\end{equation}
等价地，令 \(\mathbf{1}\) 为全 \(1\) 向量并令 \(K=\bigl(\begin{smallmatrix}1&1\\[2pt]1&0\end{smallmatrix}\bigr)\)，则有张量分解
\begin{equation}\label{eq:pom-replica-softcore-temperature-decomposition}
\boxed{\;
T^{(q)}_p
=p\,\mathbf{1}\mathbf{1}^{\mathsf T}+(1-p)\,K^{\otimes q}\;}
\end{equation}
（其中 \((K^{\otimes q})_{u,v}=\mathbf{1}_{\{u\wedge v=0\}}\)）。
当 \(p=\tfrac12\) 时，\eqref{eq:pom-replica-softcore-temperature-kernel} 退化为命题 \ref{prop:pom-multiplicity-composition-replica-softcore-transfer} 的软核 \(T^{(q)}\)。
\par
对任意有限简单图 \(G=(V,E)\)，定义 \(q\)-复制配分函数
\begin{equation}\label{eq:pom-replica-softcore-graph-partition-function}
\boxed{\;
Z_G^{(q)}(p)
:=\sum_{\sigma:V\to U_q}\ \prod_{\{x,y\}\in E}T^{(q)}_p\bigl(\sigma(x),\sigma(y)\bigr)\;}
\end{equation}
并定义独立集计数
\[
\mathsf{I}(H)
:=\#\Bigl\{\tau:V(H)\to\{0,1\}:\ \tau(x)\tau(y)=0\ \text{对所有 }\{x,y\}\in E(H)\Bigr\}.
\]
\end{definition}

\begin{remark}[单比特温度核的 NAND 表述]\label{rem:pom-replica-softcore-temperature-kernel-nand}
在定义 \ref{def:pom-replica-softcore-temperature-kernel-bernoulli} 中，单比特矩阵 \(K\) 满足
\[
K_{a,b}=1\ \Longleftrightarrow\ a\wedge b=0
\ \Longleftrightarrow\ (a,b)\neq(1,1),
\qquad a,b\in\{0,1\}.
\]
因此 \(K\) 与布尔 NAND 门（仅在 \((1,1)\) 处输出 \(0\)）的真值表矩阵同构；其张量幂 \(K^{\otimes q}\) 则编码 \(q\) 坐标同时满足 \(u\wedge v=0\) 的约束核。
\end{remark}

\begin{theorem}[指数温度调度下的 Perron 自由能相变律]\label{thm:pom-replica-softcore-temperature-free-energy-phase-transition}
\leanverified{paper\_pom\_replica\_softcore\_temperature\_free\_energy\_phase\_transition}
令 \(\varphi=\frac{1+\sqrt5}{2}\)，并令 \(\rho(\cdot)\) 表示谱半径。
对任意常数 \(\alpha\ge 0\)，取温度参数随阶数指数衰减
\[
p_q:=2^{-\alpha q}\qquad(q\in\NN).
\]
则极限
\[
F(\alpha):=\lim_{q\to\infty}\frac1q\log \rho\!\bigl(T_{p_q}^{(q)}\bigr)
\]
存在，且满足闭式
\[
F(\alpha)=\max\bigl\{\log\varphi,\ (1-\alpha)\log 2\bigr\}.
\]
等价地，存在临界指数
\[
\alpha_c:=1-\log_2\varphi
\]
使得当 \(\alpha<\alpha_c\) 时 \(F(\alpha)=(1-\alpha)\log 2\)，当 \(\alpha\ge\alpha_c\) 时 \(F(\alpha)=\log\varphi\)。
\end{theorem}

\begin{proof}
由 \eqref{eq:pom-replica-softcore-temperature-decomposition}，
\(
T^{(q)}_{p}=p\,\mathbf1\mathbf1^{\mathsf T}+(1-p)K^{\otimes q}
\)。
注意到 \(\mathbf1\mathbf1^{\mathsf T}\) 为秩一非负矩阵，其 Perron 根为 \(\mathbf1^{\mathsf T}\mathbf1=2^q\)，故
\[
\rho\!\bigl(p\,\mathbf1\mathbf1^{\mathsf T}\bigr)=p\,2^q.
\]
又由于张量谱乘法性与 \(\rho(K)=\varphi\)，有
\[
\rho(K^{\otimes q})=\rho(K)^q=\varphi^q,
\qquad
\rho\!\bigl((1-p)K^{\otimes q}\bigr)=(1-p)\varphi^q.
\]
并且按元素非负性有 \(T_p^{(q)}\ge p\,\mathbf1\mathbf1^{\mathsf T}\) 与 \(T_p^{(q)}\ge(1-p)K^{\otimes q}\)，故由谱半径对非负矩阵的单调性得到下界
\[
\rho(T_p^{(q)})\ \ge\ \max\bigl\{p\,2^q,\ (1-p)\varphi^q\bigr\}.
\]
另一方面，\(T_p^{(q)}\) 为对称非负矩阵，故 \(\rho(T_p^{(q)})=\|T_p^{(q)}\|_2\)。
由算子范数三角不等式得
\[
\rho(T_p^{(q)})
\le
\bigl\|p\,\mathbf1\mathbf1^{\mathsf T}\bigr\|_2+\bigl\|(1-p)K^{\otimes q}\bigr\|_2
=p\,2^q+(1-p)\varphi^q.
\]
因此对任意 \(p\in[0,1]\) 有夹逼
\[
\max\bigl\{p\,2^q,\ (1-p)\varphi^q\bigr\}
\ \le\
\rho(T_p^{(q)})
\ \le\
p\,2^q+(1-p)\varphi^q.
\]
取 \(p=p_q=2^{-\alpha q}\) 得 \(p_q\,2^q=2^{(1-\alpha)q}\)，且 \((1-p_q)\varphi^q=\varphi^q(1+o(1))\)。
由两项指数增长的最大指数主导律，
\[
\lim_{q\to\infty}\frac1q\log\bigl(2^{(1-\alpha)q}+\varphi^q\bigr)
=\max\bigl\{(1-\alpha)\log 2,\ \log\varphi\bigr\},
\]
结合夹逼即得所示闭式。
临界指数由 \((1-\alpha)\log2=\log\varphi\) 等价变形得到。证毕。
\end{proof}

\begin{theorem}[对称门核的 Perron--Landauer 极限律]\label{thm:pom-perron-landauer-limit-law}
令 \(A\in\RR_{\ge 0}^{d\times d}\) 为实对称非负矩阵，并记 \(\lambda:=\rho(A)\)。
对每个 \(q\ge 1\) 与任意温度调度 \(p_q\in[0,1]\)，在 \(\RR^{d^q}\cong(\RR^{d})^{\otimes q}\) 上定义
\begin{equation}\label{eq:pom-perron-landauer-tensor-soft-kernel}
\boxed{\;
T^{(q)}_{p_q}(A)
:=
p_q\,\mathbf 1\mathbf 1^{\mathsf T}+(1-p_q)\,A^{\otimes q}\;}
\end{equation}
其中 \(\mathbf 1\in\RR^{d^q}\) 为全 \(1\) 向量。
假设指数率极限
\[
L:=\lim_{q\to\infty}\frac{1}{q}\log p_q\ \in[-\infty,0]
\]
存在（约定 \(\log 0=-\infty\)）。
则极限
\[
F(A;\{p_q\})
:=\lim_{q\to\infty}\frac{1}{q}\log\rho\!\bigl(T^{(q)}_{p_q}(A)\bigr)
\]
存在并满足闭式
\[
F(A;\{p_q\})=\max\bigl\{\log\lambda,\ \log d + L\bigr\}.
\]
特别地，若 \(p_q=d^{-\alpha_q q}\) 且 \(\alpha_q\to\alpha\)，则
\[
F(A;\{p_q\})=\max\bigl\{\log\lambda,\ (1-\alpha)\log d\bigr\}.
\]
\end{theorem}
\begin{proof}
由于 \(A\) 对称，\(\rho(A)\) 为最大特征值，且对任意 \(q\) 有 \(\rho(A^{\otimes q})=\rho(A)^q=\lambda^q\)。
又因 \(\mathbf 1\mathbf 1^{\mathsf T}\) 为秩一非负对称矩阵，其唯一非零特征值为 \(\mathbf 1^{\mathsf T}\mathbf 1=d^q\)，故
\[
\rho\!\bigl(p_q\,\mathbf 1\mathbf 1^{\mathsf T}\bigr)=p_q d^q,
\qquad
\rho\!\bigl((1-p_q)A^{\otimes q}\bigr)=(1-p_q)\lambda^q.
\]

由逐项非负性有
\(
T^{(q)}_{p_q}(A)\ge p_q\,\mathbf 1\mathbf 1^{\mathsf T}
\)
与
\(
T^{(q)}_{p_q}(A)\ge (1-p_q)A^{\otimes q}
\)，
从而谱半径的 Perron 单调性给出下界
\[
\rho\!\bigl(T^{(q)}_{p_q}(A)\bigr)\ \ge\ \max\{p_q d^q,\ (1-p_q)\lambda^q\}.
\]
另一方面，\(T^{(q)}_{p_q}(A)\) 仍为实对称矩阵，故 \(\rho(\cdot)=\|\cdot\|_2\)。
由算子范数的三角不等式，
\[
\rho\!\bigl(T^{(q)}_{p_q}(A)\bigr)
\le
\bigl\|p_q\,\mathbf 1\mathbf 1^{\mathsf T}\bigr\|_2+\bigl\|(1-p_q)A^{\otimes q}\bigr\|_2
=p_q d^q+(1-p_q)\lambda^q.
\]
因此对每个 \(q\) 有夹逼
\[
\max\{p_q d^q,\ (1-p_q)\lambda^q\}
\ \le\
\rho\!\bigl(T^{(q)}_{p_q}(A)\bigr)
\ \le\
p_q d^q+(1-p_q)\lambda^q.
\]

将三项取对数并除以 \(q\)。
若 \(p_q=0\) 在无穷多 \(q\) 上出现，则由 \(L=\lim q^{-1}\log p_q\) 的存在性只能是 \(L=-\infty\) 且 \(p_q=0\) 对充分大 \(q\) 恒成立；此时 \(T^{(q)}_{p_q}(A)=A^{\otimes q}\) 对充分大 \(q\) 成立，极限值为 \(\log\lambda=\max\{\log\lambda,\log d+L\}\)，结论平凡。
下文可因此假设 \(p_q>0\) 对充分大 \(q\) 成立。
注意到 \(\lim_{q\to\infty}\frac1q\log(1-p_q)=0\)（因 \(\log(1-p_q)\) 有界且 \(q\to\infty\)），且
\[
\frac1q\log(p_q d^q)=\log d+\frac1q\log p_q\ \to\ \log d+L.
\]
由两项指数主导律
\(
\lim_{q\to\infty}\frac1q\log(x_q+y_q)=\max\{\lim\frac1q\log x_q,\lim\frac1q\log y_q\}
\)
（当两极限存在且 \(x_q,y_q>0\)）以及上述夹逼，即得极限存在并等于 \(\max\{\log\lambda,\log d+L\}\)。
最后一条取 \(p_q=d^{-\alpha_q q}\) 并令 \(L=-\alpha\log d\) 即得。证毕。
\end{proof}



\begin{theorem}[Bernoulli 子图的独立集矩表示]\label{thm:pom-replica-softcore-bernoulli-subgraph-moment-representation}
在定义 \ref{def:pom-replica-softcore-temperature-kernel-bernoulli} 的设定下，对任意有限简单图 \(G=(V,E)\)、任意整数 \(q\ge 1\) 与任意 \(p\in[0,1]\) 有严格恒等式
\begin{equation}\label{eq:pom-replica-softcore-bernoulli-subgraph-moment-representation}
\boxed{\;
Z_G^{(q)}(p)
=\sum_{F\subseteq E}p^{|E|-|F|}(1-p)^{|F|}\,\mathsf{I}(G_F)^{q}\;}
\end{equation}
其中 \(G_F:=(V,F)\) 为生成子图。
等价地，若 \(\mathbf{F}\subseteq E\) 为按边独立取样、以概率 \((1-p)\) 保留边的 Bernoulli 随机子图，则
\[
Z_G^{(q)}(p)=\mathbb{E}\bigl[\mathsf{I}(G_{\mathbf{F}})^{q}\bigr].
\]
\leanverified{paper\_pom\_replica\_bernoulli\_subgraph\_moment\_seeds}
\leanverified{paper\_pom\_replica\_softcore\_bernoulli\_subgraph\_moment\_representation}
\end{theorem}
\begin{proof}
对每条边 \(e=\{x,y\}\in E\) 定义
\(
D_e(\sigma):=\mathbf{1}_{\{\sigma(x)\wedge\sigma(y)=0\}}
\)。
由 \eqref{eq:pom-replica-softcore-temperature-kernel}，有
\(
T^{(q)}_p\bigl(\sigma(x),\sigma(y)\bigr)=p+(1-p)D_e(\sigma)
\)。
因此
\[
\prod_{e\in E}\bigl(p+(1-p)D_e(\sigma)\bigr)
=\sum_{F\subseteq E}p^{|E|-|F|}(1-p)^{|F|}\prod_{e\in F}D_e(\sigma),
\]
将其代回 \eqref{eq:pom-replica-softcore-graph-partition-function} 并交换求和次序得
\[
Z_G^{(q)}(p)
=\sum_{F\subseteq E}p^{|E|-|F|}(1-p)^{|F|}
\sum_{\sigma:V\to U_q}\prod_{e\in F}D_e(\sigma).
\]
对固定 \(F\)，写 \(\sigma=(\sigma^{(1)},\dots,\sigma^{(q)})\) 为 \(q\) 个坐标的二元赋值，其中 \(\sigma^{(a)}:V\to\{0,1\}\)。
注意到
\(
D_e(\sigma)=\prod_{a=1}^q\mathbf{1}_{\{\sigma^{(a)}(x)\sigma^{(a)}(y)=0\}}
\)，
从而
\[
\prod_{e\in F}D_e(\sigma)
=\prod_{a=1}^{q}\ \prod_{\{x,y\}\in F}\mathbf{1}_{\{\sigma^{(a)}(x)\sigma^{(a)}(y)=0\}}.
\]
因此坐标 \(a\) 间完全解耦，并且每个坐标的允许赋值数恰为 \(\mathsf{I}(G_F)\)。
故内层求和为 \(\mathsf{I}(G_F)^q\)，代回即得 \eqref{eq:pom-replica-softcore-bernoulli-subgraph-moment-representation}。证毕。
\end{proof}

\begin{corollary}[环幂迹的 Bernoulli 子图矩表示]\label{cor:pom-replica-softcore-cycle-trace-bernoulli-moment}
取环图 \(C_m\)（\(m\ge 1\)），则对任意整数 \(q\ge 1\) 与任意 \(p\in[0,1]\) 有
\begin{equation}\label{eq:pom-replica-softcore-cycle-trace-bernoulli-moment}
\boxed{\;
\Tr\!\bigl((T^{(q)}_p)^m\bigr)
=\sum_{F\subseteq E(C_m)}p^{m-|F|}(1-p)^{|F|}\,\mathsf{I}\bigl((C_m)_F\bigr)^{q}\;}
\end{equation}
其中 \((C_m)_F\) 为生成子图。
\leanverified{paper\_pom\_replica\_cycle\_trace\_bernoulli\_seeds}
\leanverified{paper\_pom\_replica\_cycle\_trace\_bernoulli\_package}
\end{corollary}
\begin{proof}
将定义 \ref{def:pom-replica-softcore-temperature-kernel-bernoulli} 的 \(Z_G^{(q)}(p)\) 取 \(G=C_m\)。
注意到对环图有矩阵恒等式 \(Z_{C_m}^{(q)}(p)=\Tr((T^{(q)}_p)^m)\)，再应用定理 \ref{thm:pom-replica-softcore-bernoulli-subgraph-moment-representation} 即得结论。证毕。
\end{proof}

\leanverified{paper\_pom\_replica\_softcore\_bernoulli\_stieltjes\_hankel}
\begin{proposition}[Stieltjes 矩结构与 Hankel 正半定]\label{prop:pom-replica-softcore-bernoulli-stieltjes-hankel}
固定有限简单图 \(G=(V,E)\) 与 \(p\in[0,1]\)，令 \(a_q:=Z_G^{(q)}(p)\)（\(q\ge 0\)，并约定 \(Z_G^{(0)}(p)=\sum_{F\subseteq E}p^{|E|-|F|}(1-p)^{|F|}=1\)）。
则存在有限支撑概率测度 \(\mu_{G,p}\) 使
\begin{equation}\label{eq:pom-replica-softcore-bernoulli-stieltjes}
\boxed{\;
a_q=\int_{[0,\infty)} x^{q}\,\dd\mu_{G,p}(x)\qquad(q\ge 0)\;}
\end{equation}
从而 \((a_q)_{q\ge 0}\) 为有限支撑的 Stieltjes 矩序列。
特别地，对任意 \(n\ge 0\)，Hankel 矩阵 \(\bigl[a_{i+j}\bigr]_{0\le i,j\le n}\) 正半定，因而全部主子式非负。
此外，若 \(\mu_{G,p}\) 的支撑至少含两点，则对所有整数 \(q\ge 1\) 有严格对数凸
\[
a_q^2<a_{q-1}a_{q+1}.
\]
\end{proposition}
\begin{proof}
由 \eqref{eq:pom-replica-softcore-bernoulli-subgraph-moment-representation}，
\(
a_q=\sum_{F\subseteq E}w(F)\,\mathsf{I}(G_F)^q
\)，其中 \(w(F)=p^{|E|-|F|}(1-p)^{|F|}\) 且 \(\sum_F w(F)=1\)。
令
\(
\mu_{G,p}:=\sum_{F\subseteq E}w(F)\,\delta_{\mathsf{I}(G_F)}
\)
即得 \eqref{eq:pom-replica-softcore-bernoulli-stieltjes}。
\par
Hankel 正半定性为 Gram 矩阵表述：对任意 \(c_0,\dots,c_n\in\RR\)，有
\[
\sum_{i,j=0}^{n}c_ic_j a_{i+j}
=\int\Bigl(\sum_{i=0}^{n}c_i x^i\Bigr)^2\dd\mu_{G,p}(x)\ge 0.
\]
严格对数凸由 Cauchy--Schwarz 在非退化分布下取严格不等式得到。证毕。
\end{proof}

\begin{proposition}[\texorpdfstring{$S_q$}{Sq}-对称降维的温度延拓]\label{prop:pom-replica-softcore-temperature-sq-reduction}
\leanverified{paper\_pom\_replica\_softcore\_temperature\_sq\_reduction}
固定整数 \(q\ge 1\) 与 \(p\in[0,1]\)。
矩阵 \(T_p^{(q)}\) 在坐标置换群 \(S_q\) 下不变，因此其 Perron 根位于 \(S_q\)-对称子空间。
在以 Hamming 权 \(k=|u|\in\{0,\dots,q\}\) 分层的基下，该限制矩阵为 \(A_p^{(q)}\in\RR^{(q+1)\times(q+1)}\)，其条目为
\begin{equation}\label{eq:pom-replica-softcore-temperature-Apq-def}
\boxed{\;
\bigl(A_p^{(q)}\bigr)_{k,l}
:=p\binom{q}{l}+(1-p)\binom{q-k}{l},
\qquad 0\le k,l\le q.\;}
\end{equation}
并且有严格等式
\begin{equation}\label{eq:pom-replica-softcore-temperature-sq-reduction}
\boxed{\;
\rho\!\bigl(T_p^{(q)}\bigr)=\rho\!\bigl(A_p^{(q)}\bigr).\;}
\end{equation}
\end{proposition}
\begin{proof}
对 \(u\in U_q\) 取 \(k=|u|\)，并取任意仅依赖权的正函数 \(f(u)=g(|u|)\)。
把 \((T_p^{(q)}f)(u)=\sum_{v}T_p^{(q)}(u,v)f(v)\) 按 \(l=|v|\) 分组：
当 \(|v|=l\) 时满足 \(u\wedge v=0\) 的 \(v\) 个数为 \(\binom{q-k}{l}\)，总数为 \(\binom{q}{l}\)。
故
\[
\sum_{\substack{v\in U_q\\ |v|=l}}T_p^{(q)}(u,v)
=\binom{q-k}{l}\cdot 1+\bigl(\binom{q}{l}-\binom{q-k}{l}\bigr)\cdot p
=p\binom{q}{l}+(1-p)\binom{q-k}{l},
\]
即为 \eqref{eq:pom-replica-softcore-temperature-Apq-def}。
由于 \(T_p^{(q)}\) 非负且 \(S_q\)-不变，Perron 根可由对称子空间内的正本征向量实现，故 \eqref{eq:pom-replica-softcore-temperature-sq-reduction} 成立。证毕。
\end{proof}

\begin{theorem}[例外块行列式的 \texorpdfstring{$(1-p)^q$}{(1-p)^q} 刚性]\label{thm:pom-replica-softcore-temperature-exceptional-determinant}
在命题 \ref{prop:pom-replica-softcore-temperature-sq-reduction} 的设定下，对任意整数 \(q\ge 1\) 与任意 \(p\in[0,1]\) 有
\begin{equation}\label{eq:pom-replica-softcore-temperature-exceptional-determinant}
\boxed{\;
\det\!\bigl(A_p^{(q)}\bigr)=(-1)^{\frac{q(q+1)}{2}}\,(1-p)^{q}.\;}
\end{equation}
等价地，若记 \(\mathrm{spec}(A_p^{(q)})=\{\lambda_i^{(q)}(p)\}_{i=0}^{q}\)，则
\[
\prod_{i=0}^{q}\lambda_i^{(q)}(p)=(-1)^{\frac{q(q+1)}{2}}\,(1-p)^{q}.
\]
\leanverified{paper\_pom\_replica\_exceptional\_det\_seeds}
\leanverified{paper\_pom\_replica\_softcore\_temperature\_exceptional\_determinant}
\end{theorem}
\begin{proof}
令 \(\mathbf{1}\in\RR^{q+1}\) 为全 \(1\) 向量，并记 \(b:=(\binom{q}{0},\dots,\binom{q}{q})^{\mathsf T}\)。
再令 \(C^{(q)}\in\RR^{(q+1)\times(q+1)}\) 满足 \(C^{(q)}_{k,l}=\binom{q-k}{l}\)。
由 \eqref{eq:pom-replica-softcore-temperature-Apq-def} 有分解
\[
A_p^{(q)}=p\,\mathbf{1}b^{\mathsf T}+(1-p)\,C^{(q)}.
\]
注意 \(C^{(q)}e_0=\mathbf{1}\)（因 \(C^{(q)}_{k,0}=1\)），故 \((C^{(q)})^{-1}\mathbf{1}=e_0\)；
并且 \(b^{\mathsf T}e_0=\binom{q}{0}=1\)。
对 \((1-p)C^{(q)}\) 应用秩一行列式引理得到
\[
\det(A_p^{(q)})
=(1-p)^{q+1}\det(C^{(q)})\Bigl(1+\frac{p}{1-p}b^{\mathsf T}(C^{(q)})^{-1}\mathbf{1}\Bigr)
=(1-p)^{q}\det(C^{(q)}).
\]
而 \(\det(C^{(q)})=(-1)^{\frac{q(q+1)}{2}}\) 已由 \eqref{eq:pom-replica-softcore-detC} 给出，代回即得 \eqref{eq:pom-replica-softcore-temperature-exceptional-determinant}。证毕。
\end{proof}

\begin{theorem}[温度延拓下的词迹公式与固定谱簇单项塌缩]\label{thm:pom-replica-softcore-temperature-word-trace-collapse}
\leanverified{paper\_pom\_replica\_softcore\_temperature\_word\_trace\_collapse}
沿用定理 \ref{thm:pom-replica-softcore-word-trace-power-sums} 的二字母词记号：\(\mathcal{W}_m=\{K,J_2\}^m\)、\(M_w=w_1\cdots w_m\)、\(\tau(w)=\Tr(M_w)\)。
对 \(w\in\mathcal{W}_m\) 记
\(
N_J(w):=\#\{i:\ w_i=J_2\}
\) 与 \(N_K(w):=\#\{i:\ w_i=K\}=m-N_J(w)\)。
则对任意整数 \(q\ge 1,m\ge 1\) 与任意 \(p\in[0,1]\) 有
\begin{equation}\label{eq:pom-replica-softcore-temperature-word-trace-full}
\boxed{\;
\Tr\!\bigl((T_p^{(q)})^{m}\bigr)
=\sum_{w\in\mathcal{W}_m}p^{N_J(w)}(1-p)^{N_K(w)}\,\tau(w)^{q}\;}
\end{equation}
并且若记例外块幂和为
\(
S_{m,\mathrm{exc}}(q;p):=\Tr\!\bigl((A_p^{(q)})^{m}\bigr)=\sum_{i=0}^{q}\lambda_i^{(q)}(p)^{m}
\)，则成立单项塌缩恒等式
\begin{equation}\label{eq:pom-replica-softcore-temperature-word-trace-collapse}
\boxed{\;
\Tr\!\bigl((T_p^{(q)})^{m}\bigr)-S_{m,\mathrm{exc}}(q;p)
=(1-p)^{m}\Bigl(L_m^{q}-\Omega_m(q)\Bigr)\;}
\end{equation}
其中 \(L_m=\Tr(K^m)\) 为 Lucas 数，而 \(\Omega_m(q)=\Tr(\mathrm{Sym}^q(K^m))\) 如 \eqref{eq:pom-replica-softcore-word-trace-collapse} 所定义。
\end{theorem}
\begin{proof}
由 \eqref{eq:pom-replica-softcore-temperature-decomposition} 与 Kronecker 幂的乘法性展开得
\[
(T_p^{(q)})^m
=\sum_{w\in\mathcal{W}_m}p^{N_J(w)}(1-p)^{N_K(w)}\,(M_w)^{\otimes q},
\]
再取迹并用 \(\Tr((M_w)^{\otimes q})=\Tr(M_w)^q=\tau(w)^q\) 即得 \eqref{eq:pom-replica-softcore-temperature-word-trace-full}。
\par
另一方面，\((M_w)^{\otimes q}\) 在 \(S_q\)-对称子空间的限制同构于 \(\mathrm{Sym}^q(M_w)\)，从而
\[
S_{m,\mathrm{exc}}(q;p)
=\sum_{w\in\mathcal{W}_m}p^{N_J(w)}(1-p)^{N_K(w)}\,\Tr\!\bigl(\mathrm{Sym}^q(M_w)\bigr).
\]
若 \(w\neq K^m\)，则 \(w\) 含至少一个 \(J_2\)，从而 \(\det(M_w)=0\)，并由定理 \ref{thm:pom-replica-softcore-word-trace-power-sums} 的论证得
\(
\Tr(\mathrm{Sym}^q(M_w))=\Tr(M_w)^q=\tau(w)^q
\)。
故两条加权词和逐项相消，仅余 \(w=K^m\) 一项，其权重为 \((1-p)^m\)，并分别贡献 \(L_m^q\) 与 \(\Omega_m(q)\)。
于是得到 \eqref{eq:pom-replica-softcore-temperature-word-trace-collapse}。证毕。
\end{proof}

\endinput
